% from ega4-8-fr.tex
\oldpage[IV-3]{47}
déjà que $g'_2$ est une immersion ouverte. On conclut bien que $g'$ est une
immersion ouverte.
\end{proof}

\begin{remark}[8.12.7]
\label{IV.8.12.7-fr}
Lorsque, dans (8.12.6), on suppose que $Y$ est un \emph{schéma affine}, la
démonstration par réduction au cas noethérien montre que, dans la
factorisation (8.12.6.1), les morphismes $f'$ et $u$ sont aussi des
morphismes \emph{de présentation finie} (1.6.2).
\end{remark}

\begin{corollary}[8.12.8]
\label{IV.8.12.8-fr}
Soient $Y$ un schéma quasi-compact tel qu'il existe un $\mathcal O_Y$-Module
ample (II, 4.5.3), $f:X\to Y$ un morphisme quasi-fini et quasi-projectif.
Alors il existe une factorisation de $f$ en
\[
X\xrightarrow{f'}Y'\xrightarrow{u}Y
\]
où $f'$ est une immersion ouverte et $u$ un morphisme fini.
\end{corollary}

\begin{proof}
L'hypothèse entraîne que $X$ s'identifie à un sous-$Y$-schéma quasi-compact
d'un $Y$-schéma de la forme $Z=\mathbf P_Y^r$ (II, 5.3.3). Il y a par suite
un voisinage ouvert quasi-compact $U$ de $X$ dans $Z$ tel que $X$ soit fermé
dans $U$~; comme $Z$ est un schéma, l'injection canonique $U\to Z$ est un
morphisme de présentation finie ((1.2.7) et (1.6.2)), donc le morphisme
composé $g:U\to Z\to Y$ est aussi un morphisme de présentation finie (le
fait que $\mathbf P_Y^r$ est de présentation finie sur $Y$ résultant
aussitôt de la définition (II, 4.1.1)). Soit $\mathcal J$ l'Idéal
quasi-cohérent de $\mathcal O_U$ définissant le sous-préschéma fermé $X$~;
comme $U$ est un schéma quasi-compact, $\mathcal J$ est limite inductive
filtrante de ses sous-Idéaux quasi-cohérents de type fini $\mathcal J_\lambda$
(I, 9.4.9). Si $X_\lambda$ est le sous-préschéma fermé de $U$ défini par
$\mathcal J_\lambda$, on a par suite $X=\bigcap_\lambda X_\lambda$. Pour tout
$y\in Y$, on a donc
$f^{-1}(y)=\bigcap_\lambda(X_\lambda\cap g^{-1}(y))$, et comme les ensembles
$X_\lambda\cap g^{-1}(y)$ sont fermés dans l'espace noethérien
$g^{-1}(y)$, il existe pour tout $y$ un indice $\lambda(y)$ tel que
$f^{-1}(y)=X_{\lambda(y)}\cap g^{-1}(y)$. Désignons par $E_\lambda$
l'ensemble des $y\in Y$ tels que la fibre $X_\lambda\cap g^{-1}(y)$ de la
restriction de $g$ à $X_\lambda$ soit un $k(y)$-préschéma fini. L'hypothèse
que $f$ est quasi-fini entraîne, en vertu de ce qui précède, que
$Y=\bigcup_\lambda E_\lambda$. Or, chacun des $X_\lambda$ est, par
définition, \emph{de présentation finie} sur $Y$~; il résulte donc de
(9.2.3) et (9.2.6)\footnote{Le lecteur vérifiera que les corollaires (8.12.8)
à (8.12.11) ne sont pas utilisés dans le \S~9.} que les $E_\lambda$ sont des
ensembles \emph{constructibles} dans le schéma $Y$~; comme ils forment un
ensemble filtrant croissant, il existe un indice $\lambda$ tel que
$E_\lambda=Y$ (1.9.9), et pour cet indice $\lambda$, le morphisme
$f_\lambda:X_\lambda\to Y$, restriction de $g$ à $X_\lambda$, est donc
quasi-fini. Comme il est de présentation finie et séparé, on peut lui
appliquer (8.12.6), et $f_\lambda$ se factorise donc en
\[
X_\lambda\xrightarrow{j_\lambda}Y'_\lambda
\xrightarrow{u_\lambda}Y
\]
où $j_\lambda$ est une immersion et $u_\lambda$ un morphisme fini. Comme
$X$ est un sous-préschéma fermé de $X_\lambda$, on a ainsi prouvé que $f$
possède la propriété (8.12.3, \emph{a)}), d'où le corollaire en vertu de
l'équivalence de (8.12.3, \emph{a)}) et (8.12.3, \emph{a$'$)}).
\end{proof}

